% SHARD-ID: AQM-45-SPEC-Z-ARITHMETIC-QUANTUM-FIELD
% SHARD-TITLE: The Arithmetic Quantum Field on \(\Spec\mathbb Z\)
% SHARD-SUMMARY: Builds the closed-prime residue-qudit net for \(\Spec\mathbb Z\) and separates the generic \(\mathbb Q\) sector.
% SHARD-SUMMARY: Proves that \(z\mapsto z+1\) is not a scheme automorphism of \(\Spec\mathbb Z\).
% SHARD-SUMMARY: Defines the internal integer residue-translation action and its squarefree residue-product extension.
% SHARD-KEYWORDS: Spec Z, prime qudits, generic point, incomplete tensor product, translation, squarefree residues

\section[The Arithmetic Quantum Field on Spec Z]{The Arithmetic Quantum Field on \(\Spec\mathbb Z\)}
\label{sec:spec-z-arithmetic-quantum-field}
\subsection{Status and Source Anchors}

This shard extends \texttt{CONVENTIONS.md} (u)--(y) and (ai).  Source-local
scheme facts are from \path{references/algebraic_geometry/stacks_project_algebra.tex}: lines
120--158, 250--274, 2972--2983, 3104--3160, 3251--3298, 3310--3342, and
4657--4703 cover prime ideals, residue fields, \(\Spec\), Zariski opens,
functoriality, quasi-compact opens, and generic points.  AQM-28 is the
finite-field finite-support pattern; this is the separate convention (ai)
\(\Spec\mathbb Z\) closed-prime extension.
\subsection{Lamport Dependency Skeleton}

\[
\begin{array}{rcl}
D1&:&X=\Spec\mathbb Z,\quad \eta=(0),\quad x_\ell=(\ell).\\
L1&:&\operatorname{Prime}(\mathbb Z)=\{0\}\cup\{(\ell):\ell\text{ prime}\}.\\
L2&:&\kappa(\eta)=\mathbb Q,\quad \kappa(x_\ell)=\mathbb F_\ell.\\
L3&:&\text{Nonempty opens are }\{\eta\}\cup\{x_\ell:\ell\notin S\}\text{ with }S\text{ finite}.\\
D2&:&E_{\rm cl}(U)=\bigoplus_{x_\ell\in U}\mathbb F_\ell^2,\quad \mathcal A_{\rm cl}(U)=\varinjlim_{S\Subset U_{\rm cl}}\bigotimes_{\ell\in S}M_\ell(\mathbb C).\\
T1&:&U\mapsto\mathcal A_{\rm cl}(U)\text{ is the closed-prime open-local AQF}.\\
D3&:&\mathcal H^\Omega_{\rm cl}=\ell^2(\bigoplus_\ell\mathbb F_\ell).\\
L4&:&\eta\text{ is a }\mathbb Q\text{-Weyl sector, not a finite qudit.}\\
T2&:&\operatorname{End}_{\rm Sch}(\Spec\mathbb Z)=\{\mathrm{id}\}.\\
D4&:&\alpha_n=\prod_\ell^{\rm loc}\operatorname{Ad}T_\ell(n\bmod \ell).\\
T3&:&\alpha:\mathbb Z\curvearrowright\mathcal A_{\rm cl}(X)\text{ is faithful and extends to }\prod_\ell\mathbb F_\ell.\\
T4&:&\alpha_n,\ n\ne0,\text{ is not implemented in the }|0\rangle\text{ sector.}
\end{array}
\]

\subsection{The Point Set and Topology D1--L3}
Every ideal of \(\mathbb Z\) has the form \((n)\).  If \((n)\) is prime, then
\(\mathbb Z/(n)\) is a domain by the source-local prime-ideal criterion.  This
happens exactly for \(n=0\) or for \(n=\pm \ell\) with \(\ell\) a rational
prime.  Hence
\[
  X=\Spec\mathbb Z=\{\eta\}\cup\{x_\ell:\ell\text{ prime}\},
  \qquad \eta=(0),\quad x_\ell=(\ell).
\]
The point \(\eta\) is generic because
\(\overline{\{\eta\}}=V(0)=X\), using the Stacks closure formula
\(\overline{\{\mathfrak p\}}=V(\mathfrak p)\).

\paragraph{Lemma \(L2\).}
The residue fields are
\[
  \kappa(\eta)=\mathbb Q,\qquad \kappa(x_\ell)=\mathbb F_\ell .
\]

\emph{Proof.}
By the Stacks residue-field definition,
\(\kappa(\mathfrak p)\) is the fraction field of
\(\mathbb Z/\mathfrak p\).  For \(\mathfrak p=(0)\) this is
\(\operatorname{Frac}(\mathbb Z)=\mathbb Q\).  For
\(\mathfrak p=(\ell)\) this is
\(\operatorname{Frac}(\mathbb Z/\ell\mathbb Z)=\mathbb F_\ell\).

\paragraph{Lemma \(L3\).}
The nonempty Zariski opens of \(X\) are exactly
\[
  U_S=X\setminus\{x_\ell:\ell\in S\}
     =\{\eta\}\cup\{x_\ell:\ell\notin S\},
  \qquad S\Subset\{\text{rational primes}\}.
\]

\emph{Proof.}
Closed subsets are \(V(n)\).  If \(n=0\), then \(V(n)=X\).  If \(n=\pm1\),
then \(V(n)=\emptyset\).  If \(n\ne0,\pm1\), then \(x_\ell\in V(n)\) exactly
when \(\ell\mid n\), and \(\eta\notin V(n)\).  Thus every proper nonempty
closed subset is a finite set of closed prime points.  Taking complements gives
the displayed opens.

\subsection{Closed-Prime Qudit Field D2--T1}
For a prime \(\ell\), attach
\[
  E_\ell=\mathbb F_\ell^2,\qquad
  \mathcal H_\ell=\ell^2(\mathbb F_\ell),\qquad
  \mathcal A_\ell=\operatorname{End}_{\mathbb C}(\mathcal H_\ell)\cong M_\ell(\mathbb C).
\]
Write a label \(e_\ell=(q_\ell,r_\ell)\).  The local Weyl operators are
\[
  T_\ell(q)|s\rangle=|s+q\rangle,\qquad
  R_\ell(r)|s\rangle=\exp(2\pi i\,rs/\ell)|s\rangle,\qquad
  W_\ell(q,r)=T_\ell(q)R_\ell(r).
\]
For a finite set \(S\) of primes,
\[
  E(S)=\bigoplus_{\ell\in S}\mathbb F_\ell^2,\quad
  \mathcal H(S)=\bigotimes_{\ell\in S}\mathcal H_\ell,\quad
  \mathcal A(S)=\bigotimes_{\ell\in S}\mathcal A_\ell .
\]
The Weyl commutator is
\[
  W_S(e)W_S(f)=
  \prod_{\ell\in S}\exp(2\pi i(r_\ell q'_\ell-q_\ell r'_\ell)/\ell)\,
  W_S(f)W_S(e),
\]
by the one-site calculation of AQM-28.

For an open \(U\), let \(U_{\rm cl}=U\cap\{x_\ell\}\), set
\(E_{\rm cl}(U)=\bigoplus_{x_\ell\in U_{\rm cl}}\mathbb F_\ell^2\) with finite
support, and define
\[
  \mathcal A_{\rm cl}(U)=
  \varinjlim_{S\Subset U_{\rm cl}}\mathcal A(S),
  \qquad a\mapsto a\otimes\mathbf 1\text{ for }S\subset S'.
\]

\paragraph{Theorem \(T1\).}
The assignment \(U\mapsto\mathcal A_{\rm cl}(U)\) is the closed-prime
open-local arithmetic quantum field on \(\Spec\mathbb Z\).  In particular,
\[
  \mathcal A_{\rm cl}(X)=
  \varinjlim_{S\Subset\{\ell\}}\bigotimes_{\ell\in S}M_\ell(\mathbb C),
\]
the algebraic infinite tensor product of one \(\ell\)-level qudit for each
rational prime \(\ell\).

\emph{Proof.}
This repeats AQM-28's finite-support construction on finite-residue closed
prime points, as a separate extension beyond AQM-28's finite-field scope.  If
\(U\subset V\), then \(U_{\rm cl}\subset V_{\rm cl}\), so extension by identity
gives a unital injective \(*\)-homomorphism
\(\mathcal A_{\rm cl}(U)\hookrightarrow\mathcal A_{\rm cl}(V)\).  Disjoint
finite supports commute because their Weyl commutator has no nontrivial local
factor.  If \(U=\bigcup_iU_i\), every label has finite support, so each support
point may be assigned to some \(U_i\); the same additive proof as AQM-28 shows
that \(\mathcal A_{\rm cl}(U)\) is generated by the images of the
\(\mathcal A_{\rm cl}(U_i)\).

\subsection{Incomplete Tensor Representation D3}
The algebra \(\mathcal A_{\rm cl}(X)\) is representation-independent.  To get a
Hilbert space, choose reference vectors
\(\Omega_\ell=|0\rangle\in\mathcal H_\ell\).  Let
\[
  \Gamma_{\rm fin}=\bigoplus_\ell\mathbb F_\ell
  =
  \{s=(s_\ell):s_\ell=0\text{ for all but finitely many }\ell\}.
\]
Define
\[
  \mathcal H^\Omega_{\rm cl}=\ell^2(\Gamma_{\rm fin}),
  \qquad |s\rangle=\bigotimes_\ell |s_\ell\rangle
\]
with the understood convention that all but finitely many tensor factors are
\(|0\rangle\).  A finite-support algebra \(\mathcal A(S)\) acts on the
coordinates in \(S\) and fixes the others.  These actions are compatible, so
\(\mathcal H^\Omega_{\rm cl}\) represents \(\mathcal A_{\rm cl}(X)\).
Changing the reference product state changes this Hilbert-space realization;
it is not extra geometry of \(\Spec\mathbb Z\).

\subsection{The Generic Point L4}
The generic point has residue field \(\mathbb Q\), so it is not a finite qudit.
If one includes it, the separate rational Weyl sector is
\[
  E_\eta=\mathbb Q^2,\qquad \mathcal H_\eta=\ell^2(\mathbb Q),
  \qquad \chi_\mathbb Q(a)=\exp(2\pi i a).
\]
With
\[
  T_\mathbb Q(q)|s\rangle=|s+q\rangle,\qquad
  R_\mathbb Q(r)|s\rangle=\chi_\mathbb Q(rs)|s\rangle,
\]
the same calculation gives \(R_\mathbb Q(r)T_\mathbb Q(q)
=\chi_\mathbb Q(rq)T_\mathbb Q(q)R_\mathbb Q(r)\).  The pairing is
nondegenerate: if \(q\ne0\), choose \(r=(2q)^{-1}\); if \(r\ne0\), choose
\(q=(2r)^{-1}\).

Because every nonempty open contains \(\eta\), set
\(\mathcal A_\eta=\operatorname{span}_{\mathbb C}\{T_\mathbb Q(q)R_\mathbb Q(r)\}\)
as an algebraic Weyl span; no C\(^*\)-completion, adelic replacement, or
Stone--von Neumann uniqueness is claimed here.  The generic-sector local rule is
\[
  \mathcal A_{\rm gen}(U)=
  \begin{cases}\mathbb C,&U=\emptyset,\\
  \mathcal A_\eta\otimes\mathcal A_{\rm cl}(U),&U\ne\emptyset.\end{cases}
\]
Thus the generic sector is present in every nonempty local algebra.  It is a
global generic-fibre layer, not another prime-local qudit.

\subsection[No Scheme Translation on X]{No Scheme Translation \(z\mapsto z+1\) on \(X\) T2}
\paragraph{Theorem \(T2\).}
Every scheme endomorphism of \(\Spec\mathbb Z\) is the identity.  Therefore
there is no nontrivial scheme-theoretic translation \(z\mapsto z+1\) on
\(\Spec\mathbb Z\).

\emph{Proof.}
By affine functoriality, scheme endomorphisms of \(\Spec\mathbb Z\) correspond
contravariantly to unital ring endomorphisms \(\varphi:\mathbb Z\to\mathbb Z\).
Such a \(\varphi\) satisfies \(\varphi(1)=1\), hence
\(\varphi(n)=n\varphi(1)=n\) for all \(n\in\mathbb Z\).  Thus
\(\varphi=\mathrm{id}\), and the induced scheme map is the identity.

Consequently the expression \(z\mapsto z+1\) belongs to
\(\mathbb A^1_\mathbb Z=\Spec\mathbb Z[z]\), where \(z\) is a coordinate.  It
does not define a self-map of the base scheme \(\Spec\mathbb Z\).

\subsection{Internal Integer Translation D4--T4}
There is nevertheless a canonical internal residue translation on the
closed-prime qudit algebra.  For \(n\in\mathbb Z\), define at prime \(\ell\)
\[
  U_\ell(n)=T_\ell(n\bmod \ell).
\]
For finite support \(S\), set
\[
  \alpha_n^S(a)=
  \left(\bigotimes_{\ell\in S}U_\ell(n)\right)
  a
  \left(\bigotimes_{\ell\in S}U_\ell(n)\right)^* .
\]
These maps are compatible with \(a\mapsto a\otimes\mathbf 1\), hence define an
automorphism \(\alpha_n\) of \(\mathcal A_{\rm cl}(X)\).  On one Weyl operator,
\[
  \alpha_n(W_\ell(q,r))
  =
  \exp(-2\pi i\,nr/\ell)\,W_\ell(q,r).
\]

\paragraph{Theorem \(T3\).}
The maps \(n\mapsto\alpha_n\) define a faithful action of \(\mathbb Z\) on
\(\mathcal A_{\rm cl}(X)\).  On \(\mathcal A(S)\) the action factors through
\(\mathbb Z/N_S\mathbb Z\), where \(N_S=\prod_{\ell\in S}\ell\).  Therefore it
extends to the squarefree compact residue product \(\prod_\ell\mathbb F_\ell\).
The induced \(\widehat{\mathbb Z}\)-action factors through coordinatewise
reduction \(\widehat{\mathbb Z}\to\prod_\ell\mathbb F_\ell\).

\emph{Proof.}
The group law follows from \(T_\ell(n)T_\ell(m)=T_\ell(n+m)\).  If
\(\alpha_n\) is the identity, then for every prime \(\ell\) and every
\(r\in\mathbb F_\ell\) the factor \(\exp(-2\pi i nr/\ell)\) is \(1\).  Hence
\(n\equiv0\pmod\ell\) for every \(\ell\), so \(n=0\).  On finite \(S\), only
the residues \(n\bmod\ell\) for \(\ell\in S\) matter, which is the same as
\(n\bmod N_S\) by the Chinese remainder theorem.  The same finite-coordinate
dependence gives the continuous extension to the product of residue fields.

\paragraph{Theorem \(T4\).}
For \(n\ne0\), the automorphism \(\alpha_n\) is not implemented by a unitary in
the reference-vector incomplete tensor representation
\(\mathcal H^\Omega_{\rm cl}\).

\emph{Proof.}
Let \(P_\ell(a)\) be the rank-one projection onto \(|a\rangle\).  If a unitary
\(U\) implemented \(\alpha_n\), then
\[
  U P_\ell(0) U^*=P_\ell(n\bmod\ell).
\]
Since \(P_\ell(0)\Omega=\Omega\), the vector \(U\Omega\) would satisfy
\[
  P_\ell(n\bmod\ell)U\Omega=U\Omega
\]
for every prime \(\ell\).  For \(n\ne0\), \(n\bmod\ell\ne0\) for all but
finitely many \(\ell\).  But \(\mathcal H^\Omega_{\rm cl}\) has basis indexed
by configurations that are zero at all but finitely many primes.  No nonzero
vector can be supported only on configurations with \(s_\ell=n\bmod\ell\) for
infinitely many \(\ell\).  Hence \(U\Omega=0\), impossible for a unitary.

\subsection{Speculative Directions}
The checked surprise is that the scheme has no translation, while the quantum
closed-prime field has a faithful arithmetic translation action by residues.
Finite experiments see only \(\mathbb Z/N_S\mathbb Z\), while the infinite
algebra sees the squarefree residue quotient \(\prod_\ell\mathbb F_\ell\); any
\(\widehat{\mathbb Z}\)-action factors through it.  The global translation by
\(1\) is an automorphism of the quasi-local algebra but not a unitary in the
finite-excitation vacuum sector: a large arithmetic symmetry, not a local
observable.

Source-gated questions now become visible.  Should the generic \(\mathbb Q\)
sector be replaced by an adelic self-dual sector?  Can a sourced
C\(^*\)-dynamical system with prime-indexed dimensions and \(\log\ell\) energy
weights recover Euler-product or zeta data without inserting them by hand?  Is
the nonimplementability of \(\alpha_1\) the first arithmetic superselection
sector in this programme?  These are questions, not conclusions.
